% =====================================================================
\chapter{Métodos adaptativos: AdaGrad, RMSProp y Adam}
\label{tema:5}
% =====================================================================

El momentum acelera el descenso en valles mal condicionados, pero usa la misma tasa de aprendizaje en todas las coordenadas. Los métodos adaptativos asignan a cada parámetro su propio paso, dividiendo su gradiente por una estimación de su escala típica. Este tema los presenta como preacondicionadores diagonales, estudia AdaGrad y RMSProp y analiza con detalle Adam, el optimizador sobre el que se construye la receta de preentrenamiento de los modelos de lenguaje. Se sigue~\cite[§12.7, §12.8 y §12.10]{d2l}, \cite[§8.5]{gbc} y los artículos originales~\cite{adagrad,rmsprop,adam}.

En todo el tema, las operaciones entre vectores ($g^{2}$, $\sqrt{s}$, cocientes) se entienden elemento a elemento.

\section{Preacondicionamiento diagonal}

\begin{definicion}[Descenso preacondicionado]\label{def:precond}
Sea $D\in\R^{n\times n}$ diagonal con entradas positivas. El \emph{descenso de gradiente preacondicionado} por $D$ es $\theta_{k+1}=\theta_k-\eta\,D^{-1}\nabla f(\theta_k)$.
\end{definicion}

\begin{proposicion}\label{prop:precond}
El descenso preacondicionado por $D$ sobre $f$ es el descenso de gradiente ordinario sobre $\tilde f(\varphi)=f(D^{-1/2}\varphi)$ en las variables $\varphi=D^{1/2}\theta$. Si $f(\theta)=\frac12\theta\T Q\theta$, el hessiano de $\tilde f$ es $D^{-1/2}QD^{-1/2}$.
\end{proposicion}

\begin{proof}
Por la regla de la cadena, $\nabla\tilde f(\varphi)=D^{-1/2}\nabla f(D^{-1/2}\varphi)$. Un paso de descenso sobre $\tilde f$ es $\varphi_{k+1}=\varphi_k-\eta D^{-1/2}\nabla f(\theta_k)$; multiplicando por $D^{-1/2}$, $\theta_{k+1}=\theta_k-\eta D^{-1}\nabla f(\theta_k)$. Para la cuadrática, $\tilde f(\varphi)=\frac12\varphi\T D^{-1/2}QD^{-1/2}\varphi$.
\end{proof}

Si $Q=\diag(a_1,\dots,a_n)$ y se toma $D=Q$, el hessiano transformado es la identidad y $\kappa=1$: el método converge en un paso con $\eta=1$. Un buen preacondicionador diagonal iguala las curvaturas de las coordenadas. Pero solo puede corregir el mal condicionamiento alineado con los ejes; si los autovectores de $Q$ están girados respecto de ellos, una matriz diagonal no basta (\cref{eje:5-precond}). Los métodos de este tema estiman $D$ a partir de los propios gradientes, sin calcular el hessiano.

\section{AdaGrad y RMSProp}

\begin{definicion}[AdaGrad]\label{def:adagrad}
Con $s_0=0$ y $\varepsilon>0$ pequeño, \emph{AdaGrad}~\cite{adagrad} acumula los cuadrados de los gradientes:
\[
s_t=s_{t-1}+g_t^{2},\qquad\theta_t=\theta_{t-1}-\eta\,\frac{g_t}{\sqrt{s_t}+\varepsilon}.
\]
\end{definicion}

\begin{proposicion}\label{prop:adagrad}
Si $\varepsilon=0$ y una coordenada recibe un gradiente constante $g\ne 0$, su paso en la iteración $t$ mide $\eta/\sqrt t$.
\end{proposicion}

\begin{proof}
$s_t=tg^{2}$, luego $\eta\,g/\sqrt{tg^{2}}=\eta\,\sign(g)/\sqrt t$.
\end{proof}

El paso de AdaGrad no depende de la escala del gradiente y decrece como $1/\sqrt t$, un calendario incorporado que es adecuado en problemas convexos pero demasiado agresivo en el entrenamiento de redes profundas, donde la acumulación sin olvido acaba paralizando el método~\cite[§8.5.1]{gbc}.

\begin{definicion}[RMSProp]\label{def:rmsprop}
\emph{RMSProp}~\cite{rmsprop} sustituye la suma por una media móvil exponencial de coeficiente $\rho\in(0,1)$:
\[
s_t=\rho\,s_{t-1}+(1-\rho)\,g_t^{2},\qquad\theta_t=\theta_{t-1}-\eta\,\frac{g_t}{\sqrt{s_t}+\varepsilon}.
\]
\end{definicion}

$\sqrt{s_t}$ estima la raíz de la media cuadrática reciente del gradiente de cada coordenada. Una coordenada con gradientes grandes recibe pasos pequeños, y viceversa. Por la \cref{prop:ema}, con gradiente constante $s_t=(1-\rho^{t})g^{2}$, de modo que el primer paso mide $\eta/\sqrt{1-\rho}$: unas $3{,}16$ veces $\eta$ con $\rho=0{,}9$. Este sesgo inicial, producido por el arranque en $s_0=0$, es lo que corrige Adam.

\section{Adam}

\begin{definicion}[Adam]\label{def:adam}
Sean $\eta>0$, $\beta_1,\beta_2\in[0,1)$, $\varepsilon>0$ y $m_0=v_0=0$. Para $t=1,2,\dots$, con $g_t$ el gradiente (estocástico) en $\theta_{t-1}$, \emph{Adam}~\cite[algoritmo~1]{adam} realiza
\begin{align*}
m_t&=\beta_1m_{t-1}+(1-\beta_1)\,g_t, & \hat m_t&=\frac{m_t}{1-\beta_1^{t}},\\
v_t&=\beta_2v_{t-1}+(1-\beta_2)\,g_t^{2}, & \hat v_t&=\frac{v_t}{1-\beta_2^{t}},\\
\theta_t&=\theta_{t-1}-\eta\,\frac{\hat m_t}{\sqrt{\hat v_t}+\varepsilon}.
\end{align*}
Los valores recomendados en el artículo original son $\eta=10^{-3}$, $\beta_1=0{,}9$, $\beta_2=0{,}999$ y $\varepsilon=10^{-8}$.
\end{definicion}

Adam combina el momentum (primer momento $m_t$, una media móvil del gradiente) con la normalización de RMSProp (segundo momento $v_t$, una media móvil de su cuadrado). Las divisiones por $1-\beta_1^{t}$ y $1-\beta_2^{t}$ son la \emph{corrección de sesgo}~\cite[§3]{adam}.

\begin{proposicion}[Corrección de sesgo]\label{prop:sesgo}
Si $g_t=g$ para todo $t$, entonces $m_t=(1-\beta_1^{t})\,g$ y $\hat m_t=g$. Más en general, si los $g_t$ son aleatorios con $\E[g_t]=\bar g$ para todo $t$, entonces $\E[m_t]=(1-\beta_1^{t})\,\bar g$ y $\E[\hat m_t]=\bar g$. Los mismos resultados valen para $v_t$ y $\hat v_t$ con $g_t^{2}$ y $\beta_2$.
\end{proposicion}

\begin{proof}
Por la \cref{prop:ema}, $m_t=(1-\beta_1)\sum_{s=1}^{t}\beta_1^{t-s}g_s$. Por linealidad de la esperanza,
\[
\E[m_t]=(1-\beta_1)\sum_{s=1}^{t}\beta_1^{t-s}\,\bar g=(1-\beta_1)\cdot\frac{1-\beta_1^{t}}{1-\beta_1}\,\bar g=(1-\beta_1^{t})\,\bar g,
\]
y dividiendo por $1-\beta_1^{t}$ se obtiene $\E[\hat m_t]=\bar g$. El caso constante es el caso particular determinista. El argumento para $v_t$ es idéntico.
\end{proof}

Sin corrección, $m_t$ y $v_t$ estarían sesgados hacia~0 al principio. Con $\beta_1=0{,}9$, $1-\beta_1^{t}$ vale $0{,}1$ en $t=1$, $0{,}19$ en $t=2$ y $0{,}6513$ en $t=10$; con $\beta_2=0{,}999$, $1-\beta_2^{t}$ vale $0{,}001$ en $t=1$ y solo $0{,}6323$ en $t=1000$: el sesgo de $v_t$ persiste durante unos miles de pasos.

\begin{proposicion}[Primer paso e invariancia de escala]\label{prop:escala}
Sea $\varepsilon=0$. (i) En $t=1$, cada coordenada con $g_1\ne 0$ se desplaza exactamente $\eta\,\sign(g_1)$. (ii) Si todos los gradientes se multiplican por una constante $c>0$, los iterados de Adam no cambian.
\end{proposicion}

\begin{proof}
(i) $\hat m_1=g_1$ y $\hat v_1=g_1^{2}$, luego $\hat m_1/\sqrt{\hat v_1}=g_1/\abs{g_1}$. (ii) Se propone en el \cref{eje:5-invariancia}.
\end{proof}

La invariancia de escala significa que $\eta$ fija directamente el tamaño del paso en el espacio de parámetros: Kingma y Ba observan que el desplazamiento de cada coordenada queda aproximadamente acotado por $\eta$~\cite[§2.1]{adam}. Esto hace a Adam mucho menos sensible que el SGD a la escala de la pérdida y a las diferencias de escala entre capas. Con $\varepsilon>0$ la invariancia es solo aproximada; $\varepsilon$ evita divisiones por cero cuando $\hat v_t\approx 0$.

\begin{ejemplo}\label{ej:adam}
Sea $\eta=0{,}01$, $\theta_0=1$, $\beta_1=0{,}9$, $\beta_2=0{,}999$, $\varepsilon$ despreciable y gradientes $g_1=0{,}5$ y $g_2=-0{,}1$:
\begin{center}\small
\begin{tabular}{@{}ccccccc@{}}
\toprule
$t$ & $m_t$ & $v_t$ & $\hat m_t$ & $\hat v_t$ & $\eta\,\hat m_t/\sqrt{\hat v_t}$ & $\theta_t$ \\
\midrule
1 & $0{,}05$ & $0{,}00025$ & $0{,}5$ & $0{,}25$ & $0{,}01$ & $0{,}99$ \\
2 & $0{,}035$ & $0{,}00025975$ & $0{,}1842$ & $0{,}1299$ & $0{,}005110$ & $0{,}98489$ \\
\bottomrule
\end{tabular}
\end{center}
El primer paso mide exactamente $\eta$. En el segundo, aunque el gradiente ha cambiado de signo, el primer momento conserva la inercia del primero y $\theta$ sigue disminuyendo, con un paso menor. Sin corrección de sesgo, el primer paso sería $\eta\cdot 0{,}05/\sqrt{0{,}00025}=\eta\cdot 0{,}1/\sqrt{0{,}001}\approx 3{,}162\,\eta$.
\end{ejemplo}

\begin{figure}[t]
\centering
\begin{tikzpicture}[node distance=5mm and 6mm,font=\small,
  caja/.style={draw,rounded corners=2pt,align=center,minimum height=9mm,text width=44mm,inner sep=2pt},
  flecha/.style={-{Stealth[length=2mm]}}]
  \node[caja] (gd) {Descenso de gradiente\\(Cauchy, 1847)};
  \node[caja,below=of gd] (sgd) {SGD con minibatch\\(Robbins--Monro, 1951)};
  \node[caja,below left=of sgd] (mom) {Momentum\\(Polyak, 1964)};
  \node[caja,below right=of sgd] (ada) {AdaGrad (2011)\\RMSProp (2012)};
  \node[caja,below=17mm of sgd] (adam) {Adam (2015): $m$, $v$\\y corrección de sesgo};
  \node[caja,below=of adam] (adamw) {AdamW (2019): decaimiento\\de pesos desacoplado};
  \draw[flecha] (gd) -- (sgd);
  \draw[flecha] (sgd) -- (mom);
  \draw[flecha] (sgd) -- (ada);
  \draw[flecha] (mom) -- (adam);
  \draw[flecha] (ada) -- (adam);
  \draw[flecha] (adam) -- (adamw);
\end{tikzpicture}
\caption{Genealogía de los optimizadores del módulo. AdamW, que se estudia en el \cref{tema:6}, es el optimizador habitual en el preentrenamiento de los modelos de lenguaje abiertos documentados.}
\label{fig:genealogia}
\end{figure}

\begin{observacion}[Adam en los modelos de lenguaje]
Los artículos de modelos de lenguaje abiertos documentan variantes de Adam con $\beta_2$ menor que el valor recomendado: LLaMA usa AdamW con $\beta_1=0{,}9$ y $\beta_2=0{,}95$~\cite[§2.3]{llama}. Con $\beta_2=0{,}95$, el segundo momento promedia efectivamente unos $1/(1-\beta_2)=20$ pasos en lugar de~1000, y se adapta antes a cambios bruscos en la escala del gradiente. El precio de Adam es la memoria: guarda $m$ y $v$, dos números adicionales por parámetro. En un modelo de $7\cdot 10^{9}$ parámetros son $1{,}4\cdot 10^{10}$ valores, 56~GB en precisión simple (módulo~17). Los optimizadores usados en Claude, Gemini y Grok no se han publicado.
\end{observacion}

\begin{lecturas}
\item Zhang \emph{et al.}~\cite{d2l}, §12.7 «Adagrad», §12.8 «RMSProp» y §12.10 «Adam».
\item Goodfellow, Bengio y Courville~\cite{gbc}, §8.5 «Algorithms with Adaptive Learning Rates» (§8.5.1--8.5.4).
\item Kingma y Ba~\cite{adam}, §2 «Algorithm» y §3 «Initialization Bias Correction».
\item Duchi, Hazan y Singer~\cite{adagrad}, artículo original de AdaGrad.
\end{lecturas}

\section*{Ejercicios del tema 5}
\addcontentsline{toc}{section}{Ejercicios del tema 5}

\begin{ejercicio}{1}{eje:5-escala-test}
Justifique si es verdadera o falsa la afirmación: «si se multiplica la pérdida por 1000, el primer paso de Adam se multiplica por 1000». Compare con lo que ocurre en el SGD.
\end{ejercicio}

\begin{ejercicio}{1}{eje:5-rmsprop}
Calcule el primer paso de RMSProp, sin corrección de sesgo y con $\varepsilon$ despreciable, para $\rho=0{,}9$ y para $\rho=0{,}999$, en función de $\eta$. Explique por qué este comportamiento motiva la corrección de sesgo de Adam.
\end{ejercicio}

\begin{ejercicio}{2}{eje:5-adam}
Aplique Adam con $\eta=0{,}01$, $\theta_0=1$, $\beta_1=0{,}9$, $\beta_2=0{,}999$ y $\varepsilon$ despreciable a los gradientes $g_1=0{,}5$ y $g_2=-0{,}1$. Calcule $\theta_1$ y $\theta_2$.
\end{ejercicio}

\begin{ejercicio}{2}{eje:5-precond}
Sea $f(x)=0{,}1x_1^{2}+2x_2^{2}$. (a) Calcule el número de condición de su hessiano y el del hessiano preacondicionado por $D=\diag(\nabla^{2}f)$. (b) Se gira el problema $45^{\circ}$: $g(y)=f(R\T y)$ con $R$ la rotación de ángulo $\pi/4$. Calcule el hessiano de $g$, su diagonal y el número de condición tras preacondicionar por esa diagonal. (c) ¿Qué concluye sobre los métodos adaptativos diagonales?
\fuente{adaptado de Zhang \emph{et al.}~\cite{d2l}, §12.7, ejercicios~2 y~4.}
\end{ejercicio}

\begin{ejercicio}{2}{eje:5-memoria}
Un modelo tiene $7\cdot 10^{9}$ parámetros y se entrena con Adam guardando parámetros, gradientes y los dos momentos en precisión simple (4 bytes por número). Calcule la memoria total y la fracción que corresponde al estado del optimizador.
\end{ejercicio}

\begin{ejercicio}{3}{eje:5-invariancia}
Demuestre la parte (ii) de la \cref{prop:escala}: con $\varepsilon=0$, si $g_t$ se sustituye por $cg_t$ con $c>0$ para todo $t$, los iterados $\theta_t$ de Adam no cambian. ¿Es cierto también para $c<0$? ¿Y para el SGD con momentum?
\end{ejercicio}

\begin{ejercicio}{3}{eje:5-sin-correccion}
Demuestre que, con $\varepsilon\to 0$, Adam con corrección de sesgo equivale a usar $m_t$ y $v_t$ sin corregir con la tasa de aprendizaje $\eta_t=\eta\sqrt{1-\beta_2^{t}}/(1-\beta_1^{t})$. Calcule $\eta_t/\eta$ para $t=1$, 2, 10 y 1000 con $\beta_1=0{,}9$ y $\beta_2=0{,}999$, y reproduzca el paso de $t=2$ del \cref{eje:5-adam}.
\fuente{adaptado de Zhang \emph{et al.}~\cite{d2l}, §12.10, ejercicio~2.}
\end{ejercicio}
